\section{Cartier indices and canonical loci}\label{sec:consequences}

The two detection theorems now give the index and locus statements. These deductions also make their different scopes explicit.

\begin{proof}[Proof of Corollary~\ref{cor:indices}]
Set $\alpha=[D]\in\Cl(\OO_{X,x})$. The hypothesis that $D$ is $\Q$-Cartier at $x$ means that $\alpha$ has finite order, equal to $\ind_x(D)$. Theorem~\ref{thm:maintorsion} makes its order equal to the order of $\rho_x(\alpha)$. In a finite direct sum the order of an element is the least common multiple of its component orders. This proves the index formula, including the empty case.

Schubert varieties are Cohen--Macaulay, so the canonical sheaf is invertible precisely at the Gorenstein points \cite[Chapter~3]{BH93}\cite[Tag 0DW9]{The26}. Apply the index formula to $K_X$ at a $\Q$-Gorenstein point. Both $\NonQGor(X)$ and $\Bbad(X)$ are contained in $\NGor(X)$. Outside their union the local canonical class is torsion and all its detector images vanish; torsion detection makes it zero. This proves the asserted equality of canonical loci.
\end{proof}

\begin{corollary}\label{cor:typeA-locus}
For type~A Schubert varieties, $\NGor(X)=\Bbad(X)$ without a $\Q$-Gorenstein hypothesis. More generally, \eqref{eq:noncartier} holds for every Weil divisor.
\end{corollary}
\begin{proof}
If $D$ is non-Cartier at $\eta_Z$, it is non-Cartier at every specialization in $Z$, because the free locus is open. Conversely, if $D$ is Cartier at all maximal singular generic points through $x$, Theorem~\ref{thm:mainA} makes its local class at $x$ zero. Taking $D=K_X$ proves the canonical assertion.
\end{proof}

\begin{corollary}[The factorial locus in type A]\label{cor:factorial}
For a type~A Schubert variety,
\[
 \{x:\OO_{X,x}\text{ is not factorial}\}
 =\bigcup_{\substack{Z\in\M(X)\\\OO_{X,\eta_Z}\text{ is not factorial}}}Z.
\]
\end{corollary}
\begin{proof}
If all target class groups through $x$ vanish, injectivity of $\rho_x$ makes $\Cl(\OO_{X,x})$ vanish. A normal Noetherian local domain is factorial exactly when its class group is zero \cite{Fos73}. Conversely, if $\OO_{X,x}$ is factorial, each localization $\OO_{X,\eta_Z}$ is factorial.
\end{proof}

This is the geometric locus conclusion in \cite[Conjecture~6.14]{WY08}\cite[Conjecture~7.39]{WY23}. The global permutation-avoidance classification in \cite{BMB07} addresses a different assertion.

\begin{remark}\label{rem:source-torsion}
The index formula assumes that the local class $[D]\in\Cl(\OO_{X,x})$ is torsion. It neither tests $\Q$-Cartierness from the target indices nor assigns a finite index to an infinite-order class. In particular, generic indices equal to one do not remove $\NonQGor(X)$ from the identity in Corollary~\ref{cor:indices}. The next section exhibits that obstruction.
\end{remark}
